\section{Trace-conditioned opportunity}
\label{sec:trace-results}

\claimid{RC-007}All \TraceReplayRepetitions{} cell-seed rows from nine generated swarms pass
the five validity gates. In
the prospectively frozen primary cell, the replay generates a mean of
\TracePrimaryEvents{} events. Fixed windows admit \TraceFixedPct\%, while exact
sliding-deadline packing admits \TraceExactPct\%. The exact share ranges from
\TracePrimaryExactMinPct\% to \TracePrimaryExactMaxPct\% across the three
seeds. The local bound is \TraceUpperPct\%.

\begin{table}[t]
  \centering
  \footnotesize
  % Generated by scripts/build_paper_artifacts.py. Do not edit.
\begin{tabular}{@{}lr@{}}
  \toprule
  Quantity & Primary value \\
  \midrule
  Fixed-partition share $F$ & 30.19\% \\
  Exact offline share $P^\star$ & 43.00\% \\
  Local-overlap bound $U$ & 45.85\% \\
  Alignment-gap closure & 81.83\% \\
  Mean generated events & 651{,}123 \\
  Mean exact batches & 1{,}046.3 \\
  \bottomrule
\end{tabular}

  \caption{Prospectively frozen primary trace cell: 100{,}000 target active sessions,
  route-key grouping, $K=256$, and a 50 ms launch deadline. Shares and counts are
  means over three replay seeds conditional on one fixed panel.}
  \label{tab:trace-primary}
\end{table}

Exact packing therefore gains \TraceFixedToExactPP{} percentage points over the
frozen partition and remains \TraceExactToUpperPP{} points below the local
upper bound. The mean of the per-seed gap closures is
\TraceGapClosurePct\%. This distinction matters operationally: the fixed
partition leaves model-admissible cross-boundary cohorts unused, while $U$ still
overcounts opportunities that cannot all be selected together.

\begin{figure*}[t]
  \centering
  \includegraphics[width=0.82\linewidth]{trace-exact-opportunity-frontier.png}
  \caption{Route-key-conditioned opportunity at $C=100{,}000$ and $K=256$. The
  prospectively frozen primary deadline is 50 ms. Values are means across three replay
  seeds conditional on one panel and one arrival model, not population
  confidence intervals.}
  \label{fig:trace-frontier}
\end{figure*}

\subsection{Cohort supply collapses below the boundary}

Under route-key grouping at $K=256$, $\pstar$ is zero for every tested deadline when
$C\leq10{,}000$. Even at $C=100{,}000$, it remains zero at 10 and 25 ms. Thus
a large nominal swarm does not imply that a route sees a profitable cohort
inside a short launch budget.

Reducing $K$ changes the boundary sharply. At $C=10{,}000$ and 50 ms,
$\pstar$ is \TracePAtTenKThirtyTwoPct\% for $K=32$ but
\TracePAtTenKSixtyFourPct\% for $K=64$. At $C=100{,}000$ and 50 ms, the exact
shares for $K=32,64,128,256$ are respectively
\TracePAtHundredKThirtyTwoPct\%, \TracePAtHundredKSixtyFourPct\%,
\TracePAtHundredKOneTwentyEightPct\%, and
\TracePAtHundredKTwoFiftySixPct\%. Lowering the hardware threshold can matter
more than admitting another boundary-crossing cohort.

The surface gives an online runtime a measurable workload budget. Under the
model, $F$ is the fixed-window baseline, $\pstar$ is the most that any legal
scheduler can attain, and $\pstar-F$ is the headroom available to a better
packing policy. A real implementation must measure its route-specific $K$ and
report how much of that headroom becomes achieved share $A$.
